\documentclass[14pt,a4paper]{extarticle}
\usepackage{amsmath}
\usepackage{graphicx}
\usepackage[utf8]{inputenc}    
\usepackage[T2A]{fontenc}
\usepackage[russian]{babel}
\usepackage{setspace}
\usepackage{geometry}
\usepackage{longtable}
\usepackage{makecell}
\usepackage{listings}
\usepackage{xcolor} 
\usepackage{siunitx}


\geometry{
    left=3cm,
    right=1.5cm,
    top=2cm,
    bottom=2cm
}

\begin{document}
\thispagestyle{empty}

\begin{center}
САНКТ-ПЕТЕРБУРГСКИЙ ГОСУДАРСТВЕННЫЙ УНИВЕРСИТЕТ\\
МАТЕМАТИКО-МЕХАНИЧЕСКИЙ ФАКУЛЬТЕТ\\
КАФЕДРА ФИЗИЧЕСКОЙ МЕХАНИКИ
\end{center}

\vspace{4cm}

\begin{center}
\textbf{Методы измерений и электромеханические системы}\\
\vspace{0.5cm}
Отчёт по лабораторной работе №4
\end{center}

\vspace{0.5cm}

\begin{center}
\textbf{«Измерение ускорения свободного падения с помощью электронного хронометра Ч3-32»}
\end{center}

\vspace{4cm}

\begin{flushright}
Выполнил студент:\\
Тузова Анастасия Александровна\\
группа: 24.Б51-мм

\vspace{1.5cm}

Проверил:\\
к.ф.-м.н., доцент\\
Кац Виктор Михайлович
\end{flushright}

\vfill

\begin{center}
Санкт-Петербург, 2026 г.
\end{center}
\newpage  
\tableofcontents 
\newpage
\section{Введение}
\subsection{Цель работы}
 
\begin{enumerate}
    \item Проверка принципа эквивалентности масс.
    \item Измерение ускорения свободного падения тел.
    \item Знакомство с методом измерения интервалов времени между импульсами частотомером – хронометром Ч3-32.
    \item Определение погрешности косвенных измерений.
\end{enumerate}
\newpage
\section{Основная часть}
\subsection{Теоретическая часть}
Массу тела можно определить путем измерения испытываемого телом ускорения \(\vec{a}\) под действием известной силы \(\vec{F}\):
\[
M_{\text{ин}} = \frac{|\vec{F}|}{|\vec{a}|}.
\]
Определяемая таким путем масса известна под названием инертной массы. Массу тела можно также определить, измеряя силу тяготения к другому телу, например, к Земле
\[
F = \gamma \frac{M_{\text{гp}} M_3}{r^2}; \qquad 
M_{\text{гp}} = \frac{F r^2}{\gamma M_3};
\]
здесь \(M_3\) — масса Земли,
\(\gamma\) — гравитационная постоянная,
\(r\) — расстояние между центром Земли и телом.

Определяемая таким образом масса \(M_{\text{cp}}\) называется гравитационной массой.
\\Поскольку из закона Галилея следует, что \( a_1 = a_2 = g \), то
\[
\frac{M_{\text{ин}}^{(1)}}{M_{\text{ин}}^{(2)}} = \frac{M_{\text{гр}}^{(1)}}{M_{\text{гр}}^{(2)}}, \qquad 
\frac{M_{\text{ин}}^{(1)}}{M_{\text{гр}}^{(1)}} = \frac{M_{\text{ин}}^{(2)}}{M_{\text{гр}}^{(2)}} = \text{const}. \tag{3}
\]
Коль скоро это отношение постоянно для всех тел, можно всегда привести его к единице, подобрав подходящее значение к гравитационной постоянной \( \gamma \).
\\Закон эквивалентности масс:
\[
M_{\text{ин}} = M_{\text{гр}}
\]
\\Масса шарика вычисляется по формуле:
\[
m = \rho \cdot V,
\]
где объем шара \(V\) равен
\[
V = \frac{4}{3}\pi R^3 = \frac{\pi d^3}{6},
\]
так как радиус \(R = d/2\) (половина диаметра).
\\Уравнение движения при свободном падении имеет вид:
\[
h = v_0 t + \frac{g t^2}{2},
\]
где \(v_0\) — скорость шарика при подходе к верхнему лучу (указана на установке). Следовательно,
\[
g = \frac{2(h - v_0 t)}{t^2}.
\]
\newpage
\subsection{Эксперимент}
\begin{figure}[h!]
    \centering
   \includegraphics[width=0.8\textwidth, trim=0cm 8cm 0cm 0cm, clip]{Notes_260314_003524.pdf}
    \caption{Схема установки}
    \label{fig:img1}
\end{figure}
\newpage
\subsection{Обработка данных и обсуждение результатов}
\subsubsection*{Исходный код}
Рассчитаем массу по формуле V шара и данной плотности 
\begin{lstlisting}[language=C++]
#include <iostream>
int main() {
	double a[6];
	for (int i = 0; i < 6; ++i) {
		std::cin >> a[i];
	}
	for (int i = 0; i < 6; ++i) {
		std::cout << (1.0 * a[i] * 3.1415926) / 6 
        << std::endl;
	}
	return 0;
}
\end{lstlisting}
\begin{lstlisting}[language=C++]
#include <iostream>
#include<cmath>
int main() {
	double a[30];
	double aver;
	std::cin >> aver;
	for (int i = 0; i < 30; ++i) {
		std::cin >> a[i];
	}
	double sum = 0;
	for (int i = 0; i < 30; ++i) {
		sum += (a[i] - aver)*(a[i] - aver);
	}
	std::cout << sqrt(sum / (30 * 29));
	return 0;
}
\end{lstlisting}
\newpage
\subsubsection*{Таблицы}
\begin{table}[h!]
\small
\centering
\caption{Массы шариков}
\label{tab:1}
\begin{tabular}{|c|l|c|c|c|}
\hline
№№ п/п & Вещество & Плотность, \(10^3\,\text{кг/м}^3\) & Диаметр, \(10^{-3}\,\text{м}\) & Масса, \(10^{-6}\,\text{кг}\) \\
\hline
1 & дерево (берёза) & 0,70 & 1,00 & 0,3665 \\
2 & плексиглас      & 1,18 & 1,00 & 0,6178 \\
3 & алюминий        & 2,79 & 1,00 & 1,4608 \\
4 & сталь           & 7,90 & 1,00 & 4,1364 \\
5 & латунь          & 8,50 & 1,00 & 4,4506 \\
6 & свинец          & 11,34 & 1,00 & 5,9376 \\
\hline
\end{tabular}
\end{table}
\begin{table}[htbp]
\centering
\caption{Время падения шариков}
\label{tab:falling_times}
\footnotesize  % или \scriptsize
\begin{tabular}{|c|S[table-format=3.3]|S[table-format=3.3]|S[table-format=3.3]|S[table-format=3.3]|S[table-format=3.3]|S[table-format=3.3]|}
\hline
№ & Алюминий & Латунь & Сталь & Дерево & Плексиглас & Свинец \\
\hline
1 & 152.221 & 152.047 & 152.202 & 153.129 & 153.853 & 153.408 \\
2 & 153.287 & 151.897 & 151.820 & 153.915 & 153.236 & 152.424 \\
3 & 152.563 & 151.962 & 151.940 & 154.291 & 153.092 & 152.137 \\
4 & 151.571 & 152.239 & 151.783 & 153.208 & 152.977 & 152.558 \\
5 & 151.628 & 151.737 & 151.979 & 152.950 & 153.051 & 152.241 \\
6 & 151.859 & 152.866 & 152.012 & 153.272 & 153.030 & 152.776 \\
7 & 152.180 & 151.969 & 151.324 & 153.582 & 153.346 & 152.589 \\
8 & 152.884 & 151.961 & 152.861 & 153.471 & 152.936 & 152.717 \\
9 & 152.309 & 151.920 & 151.932 & 152.995 & 153.051 & 153.070 \\
10 & 152.469 & 151.619 & 152.060 & 153.059 & 152.927 & 153.447 \\
11 & 152.264 & 151.828 & 152.891 & 153.617 & 153.146 & 152.937 \\
12 & 151.871 & 151.470 & 152.149 & 153.024 & 153.156 & 153.117 \\
13 & 151.658 & 152.029 & 152.860 & 153.016 & 153.739 & 152.902 \\
14 & 152.103 & 152.071 & 152.214 & 153.065 & 153.385 & 153.596 \\
15 & 151.862 & 151.806 & 153.222 & 153.238 & 152.977 & 152.411 \\
16 & 151.638 & 151.571 & 152.007 & 153.750 & 152.914 & 152.276 \\
17 & 152.018 & 151.433 & 152.124 & 153.249 & 152.882 & 153.144 \\
18 & 151.825 & 152.181 & 152.155 & 153.737 & 153.134 & 152.995 \\
19 & 152.003 & 151.662 & 151.838 & 152.988 & 153.279 & 152.972 \\
20 & 152.247 & 151.585 & 152.160 & 153.507 & 153.206 & 152.615 \\
21 & 152.279 & 151.586 & 152.076 & 154.145 & 153.069 & 152.511 \\
22 & 152.996 & 152.338 & 151.996 & 154.589 & 153.096 & 152.390 \\
23 & 153.855 & 152.067 & 152.151 & 154.004 & 153.097 & 152.583 \\
24 & 151.750 & 151.948 & 152.027 & 153.819 & 153.219 & 152.929 \\
25 & 152.254 & 151.875 & 152.603 & 154.280 & 152.822 & 152.879 \\
26 & 152.273 & 152.001 & 153.121 & 153.503 & 153.001 & 152.587 \\
27 & 151.550 & 152.381 & 152.632 & 153.481 & 152.886 & 152.423 \\
28 & 152.086 & 152.714 & 151.975 & 153.602 & 152.919 & 153.531 \\
29 & 152.104 & 151.935 & 152.031 & 154.747 & 154.144 & 152.633 \\
30 & 151.983 & 151.927 & 152.256 & 153.900 & 153.080 & 153.288 \\
\hline
\multicolumn{1}{|c|}{Среднее} & 152.186 & 151.954 & 152.213 & 153.571 & 153.155 & 152.803 \\
\hline
\end{tabular}
\end{table}
\newpage
Среднее значение:
\[
\bar{t} = \frac{1}{n} \sum_{i=1}^n t_i
\]

Стандартное отклонение (абсолютная погрешность):
\[
\Delta t = \sqrt{\frac{\sum_{i=1}^n (t_i - \bar{t})^2}{n(n-1)}}
\]
\begin{table}[h!]
\centering
\caption{Погрешность времени пролёта шариков \(\Delta t\)}
\label{tab:delta_t_wide}
\renewcommand{\arraystretch}{1.3}  
\setlength{\tabcolsep}{10pt}   

\begin{tabular}{|>{\centering\arraybackslash}m{3cm}|>{\centering\arraybackslash}m{3cm}|}

\hline
 \textbf{Вещество} & \textbf{\(\Delta t\), мс} \\
\hline
Алюминий    & 0.095 \\
\hline
Латунь      & 0.060 \\
\hline
Сталь       & 0.078 \\
\hline
Дерево      & 0.0906 \\
\hline
Плексиглас  & 0.054 \\
\hline
Свинец      & 0.0726 \\
\hline
\end{tabular}
\end{table}
\begin{table}[h!]
\centering
\caption{Ускорение свободного падения для различных материалов}
\label{tab:g_values}
\begin{tabular}{|c|l|c|}
\hline
№ п/п & Вещество & \( g, \, \text{м/с}^2 \) \\
\hline
1 & Алюминий    &9.68928 \\
\hline
2 & Латунь      &9.73999 \\
\hline
3 & Сталь       &9.6834 \\
\hline
4 & Дерево      & 9.39198\\
\hline
5 & Плексиглас  &9.48031 \\
\hline
6 & Свинец      & 9.5557\\
\hline
\end{tabular}
\end{table}
Формула для вычисления \(\Delta g\) (косвенные измерения):
Для функции
\[
g = \frac{2(h - v_0 t)}{t^2}
\]

Абсолютная погрешность вычисляется по формуле:
\[
\Delta g = \sqrt{\left( \frac{\partial g}{\partial h} \Delta h \right)^2 + \left( \frac{\partial g}{\partial v_0} \Delta v_0 \right)^2 + \left( \frac{\partial g}{\partial t} \Delta t \right)^2}
\]
\\Где:
\[
\Delta h = 0{,}001 \, \text{м}, \quad \Delta v_0 = 0{,}005 \, \text{м/с},
\]
\(\Delta t\) — из таблицы 3 (для каждого материала).

Частные производные:
\[
\frac{\partial g}{\partial h} = \frac{2}{t^2}
\]
\[
\frac{\partial g}{\partial v_0} = -\frac{2}{t}
\]
\[
\frac{\partial g}{\partial t} = -\frac{4(h - v_0 t)}{t^3} - \frac{2v_0}{t^2} = -\frac{4h - 2v_0 t}{t^3}
\]

Или в более удобном виде:
\[
\Delta g = g \cdot \sqrt{ \left( \frac{\Delta h}{h - v_0 t} \right)^2 + \left( \frac{\Delta v_0}{h - v_0 t} \right)^2 + \left( 2\frac{\Delta t}{t} \right)^2 }
\]

\begin{table}[h!]
\centering
\caption{Погрешность измерения ускорения свободного падения \(\Delta g\)}
\label{tab:delta_g}
\begin{tabular}{|l|c|}

\hline
\textbf{Вещество} & \textbf{\(\Delta g, \, \text{м/с}^2\)} \\
\hline
Алюминий    & 0.4405 \\
\hline
Латунь      & 0.4417\\
\hline
Сталь       & 0.4403\\
\hline
Дерево      & 0.4326\\
\hline
Плексиглас  & 0.4348\\
\hline
Свинец      & 0.4369\\
\hline

\end{tabular}
\end{table}
Самые низкие значения у легких материалов с малой плотностью - дерево, плексиглас. Это происходит из-за действующих на шарики сил сопротивления воздуха, так как они обладают меньшей инерцией и сильнее тормозятся. 

В данной работе предполагается, что движение происходит в вакууме, однако в реальных условиях сопротивление воздуха приводит к увеличению времени падения по сравнению с теоретическим (без учёта сопротивления). Завышенное значение t приводит к занижению расчётного значения g.

Также электронный хронометр и датчик имеют конечное время срабатывания, поэтому момент фиксации может немного отличаться от реального момента прохождения шарика.

Поскольку шарик имеет реальные размеры и не является материальной точкой, это может вносить дополнительную систематическую погрешность в измерение ускорения свободного падения.
\newpage
\section{Выводы}
Измерено ускорение свободного падения для шести шариков из различных материалов (алюминий, латунь, сталь, дерево, плексиглас, свинец). Полученные значения g лежат в пределах от 9.39 до 9.74 \(\text{м/с}^2\), что близко к табличному значению 9.81 \(\text{м/с}^2\)
Небольшое систематическое занижение связано с влиянием сопротивления воздуха, погрешностью прибора и размером шариков.

Определена погрешность косвенных измерений \(\Delta g\)

Для всех образцов она составила $\approx 0.44\ \text{м/с}^2$, что соответствует относительной погрешности около 4,4 \% от измеряемой величины. Столь малый разброс  \(\Delta g\) между разными материалами свидетельствует о высокой стабильности экспериментальной установки и о том, что основной вклад в погрешность вносит систематическая ошибка, одинаковая для всех измерений.
\end{document}
